\documentclass[12pt,a4paper]{article}
\usepackage[utf8]{inputenc}
\usepackage{graphicx}
\usepackage{amsmath}
\usepackage{hyperref}
\usepackage{geometry}
\usepackage{fancyhdr}
\usepackage{setspace}
\usepackage{cite}
\usepackage{url}

\geometry{margin=1in}
\onehalfspacing

\hypersetup{
    colorlinks=true,
    linkcolor=blue,
    filecolor=magenta,
    urlcolor=blue,
    pdftitle={Smart Attendance System Using Multimodal Sensors},
    pdfpagemode=FullScreen,
}

\begin{document}

% Title Page
\begin{titlepage}
	\centering
	\vspace*{2cm}

	\LARGE\textbf{Smart Attendance System Using Multimodal Sensors}\\[1cm]

	\Large\textbf{Project Report}\\[0.5cm]

	\large{Submitted in partial fulfillment of the requirements of}\\[0.3cm]
	\large{Bachelor of Technology in Electronics and Communication Engineering}\\[2cm]

	\Large\textbf{By}\\[0.5cm]
	\large{Student Name 1 (2021CS001)}\\
	\large{Student Name 2 (2021CS002)}\\
	\large{Student Name 3 (2021CS003)}\\[2cm]

	\Large\textbf{Under the Guidance of}\\[0.5cm]
	\large{Dr. Faculty Name}\\
	\large{Department of Electronics and Communication Engineering}\\[2cm]

	\Large\textbf{VIT University}\\
	\Large\textbf{Department of Electronics and Communication Engineering}\\[1cm]
	\large{\today}

	\pagebreak
\end{titlepage}

% Abstract
\begin{abstract}
	Manual attendance tracking in educational institutions is time-consuming, prone to errors, and lacks real-time monitoring capabilities. This project presents a Smart Attendance System using multimodal sensors that automates the attendance process through RFID-based teacher authentication and AI-powered facial recognition for student identification. The system utilizes an ESP32 microcontroller equipped with an RC522 RFID module for teacher authentication and an ESP32-CAM module for capturing student faces. A Flask-based web server processes the captured images using the SCRFD-34GF face detection model and ArcFace recognition algorithm. When a teacher initiates a session by scanning their RFID card, a 15-minute attendance window begins, during which the camera captures student faces. The system identifies students using facial recognition and automatically marks their attendance. Upon session completion, an automated email report is sent to the teacher containing the attendance summary. The proposed system demonstrates 85\% recognition accuracy with an average processing time of 500ms per image, offering a reliable and efficient alternative to traditional attendance methods.

	\textbf{Keywords:} Attendance System, Face Recognition, RFID, ESP32, SCRFD, ArcFace, IoT, Multimodal Sensors
\end{abstract}
\pagebreak

% Table of Contents
\tableofcontents
\pagebreak

% List of Figures
\listoffigures
\pagebreak

% List of Tables
\listoftables
\pagebreak

% Introduction
\section{Introduction}
\pagenumbering{arabic}

Attendance management is a fundamental administrative task in educational institutions. Traditional methods of attendance taking, such as calling names or physical signature sheets, consume valuable instructional time and are susceptible to proxy attendance. With the advancement of Internet of Things (IoT) technologies and artificial intelligence, automated attendance systems have become increasingly feasible and practical.

This project proposes a Smart Attendance System that combines multiple sensor modalities to create a robust and automated attendance tracking solution. The system employs RFID (Radio Frequency Identification) technology for teacher authentication and computer vision techniques for student identification. When a teacher initiates an attendance session by scanning their RFID card, a 15-minute timed window opens during which students' faces are captured and automatically recognized.

The primary objectives of this project are:
\begin{itemize}
	\item To develop a fully automated attendance system eliminating manual intervention
	\item To implement secure teacher authentication using RFID technology
	\item To achieve accurate student identification using state-of-the-art face recognition
	\item To provide real-time attendance monitoring and automated reporting
	\item To create a cost-effective solution using widely available hardware components
\end{itemize}

The proposed system addresses several limitations of existing attendance systems, including proxy attendance, time consumption, and human error. By leveraging the ESP32 microcontroller's Wi-Fi capabilities and the SCRFD-34GF face detection model, we achieve a balance between computational efficiency and recognition accuracy.

The remainder of this report is organized as follows: Section 2 presents a literature survey of existing attendance systems; Section 3 describes the methodology; Section 4 details the hardware and software components; Section 5 presents the circuit design; Section 6 discusses the results and discussion; and Section 7 concludes with future work.

\pagebreak

% Literature Survey
\section{Literature Survey}

\subsection{Overview of Attendance Systems}
Attendance tracking has evolved significantly over the decades, from paper-based methods to biometric systems. Early automated systems relied on barcode scanners and magnetic stripe cards, which required physical contact and were prone to wear and tear.

\subsection{Existing Approaches}

\subsubsection{RFID-Based Systems}
RFID technology has been widely adopted for access control and attendance tracking. Studies by \cite{RFID_attendance} have shown that RFID-based attendance systems reduce processing time by 60\% compared to manual methods. However, these systems require students to carry RFID cards and are susceptible to card sharing.

\subsubsection{Biometric Fingerprint Systems}
Fingerprint-based attendance systems, as discussed in \cite{fingerprint_system}, offer improved accuracy but face challenges with hygiene concerns, especially in post-pandemic scenarios. Additionally, fingerprint recognition can fail for individuals with worn fingerprints or certain occupations.

\subsubsection{Facial Recognition Systems}
Face recognition for attendance has gained significant attention due to its contactless nature. The work in \cite{face_recognition_survey} demonstrates that modern CNN-based face recognition systems achieve accuracy rates exceeding 99\%. However, single-modal face recognition systems can be vulnerable to photo-based spoofing attacks.

\subsection{Multimodal Approaches}
Recent research has explored combining multiple biometric modalities for improved security and accuracy. The multimodal attendance system proposed by \cite{multimodal_attendance} combines face recognition with voice verification, achieving a 40\% reduction in false acceptance rates.

\subsection{Gap Analysis}
While existing systems address individual aspects of attendance tracking, there is a need for an integrated system that combines:
\begin{itemize}
	\item Secure teacher authentication
	\item Automated student identification
	\item Real-time monitoring
	\item Automated reporting
\end{itemize}

This project bridges these gaps by implementing a multimodal approach using RFID for authentication and face recognition for identification.

\subsection{Relevant Technologies}

\subsubsection{SCRFD Face Detection}
The SCRFD (Scale-aware Face Detection) algorithm \cite{scrfd} is a state-of-the-art face detection method that achieves excellent performance across different face scales while maintaining real-time inference speed. The SCRFD-34GF variant offers the best balance of accuracy and speed for edge deployment.

\subsubsection{ArcFace Recognition}
ArcFace \cite{arcface} is a margin-based loss function for deep face recognition that enhances discriminative power of the learned embeddings. It has become the de facto standard for face recognition tasks due to its superior performance on benchmark datasets.

\pagebreak

% Methodology
\section{Methodology}

\subsection{System Overview}
The proposed Smart Attendance System operates on a client-server architecture as illustrated in Figure \ref{fig:system_architecture}. The ESP32 microcontroller serves as the edge device, collecting sensor data and communicating with the Flask-based server via Wi-Fi.

\begin{figure}[htbp]
	\centering
	\fbox{\includegraphics[width=0.8\textwidth, height=8cm]{figures/system_architecture.png}}
	\caption{System Architecture}
	\label{fig:system_architecture}
\end{figure}

\subsection{Workflow}

The system operates in the following sequence:

\subsubsection{Step 1: Teacher RFID Authentication}
The teacher initiates an attendance session by scanning their RFID card on the RC522 reader. The ESP32 reads the card's unique identifier (UID) and sends it to the server for validation. If the UID matches a registered teacher, the server creates a new attendance session and returns the session parameters.

\subsubsection{Step 2: Session Initialization}
Upon successful authentication, the server starts a 15-minute timer and configures the ESP32 to begin capturing images from the camera module. The timer ensures that attendance is captured only during the designated window.

\subsubsection{Step 3: Image Capture}
The ESP32-CAM module captures images at 5-second intervals, resulting in approximately 180 images per session. Each image is compressed as JPEG and transmitted to the server via HTTP POST requests.

\subsubsection{Step 4: Face Detection and Recognition}
Upon receiving an image, the server's face recognition module processes it through the following pipeline:
\begin{enumerate}
	\item Face detection using SCRFD-34GF
	\item Face alignment using five facial landmarks
	\item Feature extraction using ArcFace embedding network
	\item Matching against registered student database
\end{enumerate}

\subsubsection{Step 5: Attendance Recording}
Detected and recognized faces are matched to student records. The attendance is marked in the database with a confidence score. Students whose faces are detected with confidence above the threshold (0.5) are marked as present.

\subsubsection{Step 6: Session Completion}
When the 15-minute timer expires, the server automatically closes the session. An email report is generated and sent to the teacher's registered email address containing the attendance summary.

\subsection{Data Flow Diagram}

\begin{figure}[htbp]
	\centering
	\fbox{\includegraphics[width=0.8\textwidth, height=8cm]{figures/data_flow.png}}
	\caption{Data Flow Diagram}
	\label{fig:data_flow}
\end{figure}

\pagebreak

% Hardware and Software Used
\section{Hardware and Software Used}

\subsection{Hardware Components}

\subsubsection{ESP32 DevKit v1}
The ESP32 DevKit v1 is the primary microcontroller unit (MCU) for this project. Its specifications are detailed in Table \ref{tab:esp32_specs}.

\begin{table}[htbp]
	\centering
	\caption{ESP32 DevKit v1 Specifications}
	\label{tab:esp32_specs}
	\begin{tabular}{|l|c|}
		\hline
		\textbf{Parameter} & \textbf{Specification} \\
		\hline
		Processor          & Xtensa LX6 Dual-Core   \\
		Clock Speed        & 240 MHz                \\
		SRAM               & 520 KB                 \\
		Flash Memory       & 4 MB                   \\
		WiFi Standard      & 802.11 b/g/n           \\
		Bluetooth          & BLE 4.2                \\
		Operating Voltage  & 3.3V                   \\
		GPIO Pins          & 34                     \\
		\hline
	\end{tabular}
\end{table}

The ESP32's built-in Wi-Fi capability enables seamless communication with the server without additional networking hardware.

\subsubsection{RC522 RFID Module}
The MIFARE RC522 module provides contactless communication at 13.56 MHz. It communicates with the ESP32 via the SPI interface.

\begin{itemize}
	\item \textbf{Operating Frequency:} 13.56 MHz
	\item \textbf{Communication Protocol:} SPI (10 Mbps max)
	\item \textbf{Read Range:} 25mm - 80mm
	\item \textbf{Supported Cards:} MIFARE S50, S70, UltraLight
\end{itemize}

\subsubsection{ESP32-CAM Module}
The ESP32-CAM integrates the OV2640 camera sensor with the ESP32 S module.

\begin{itemize}
	\item \textbf{Sensor:} OV2640 (2 Megapixel)
	\item \textbf{Resolution:} Up to 1622x1200 pixels
	\item \textbf{JPEG Output:} Up to 1600x1200
	\item \textbf{Field of View:} 120 degrees
	\item \textbf{Flash LED:} Built-in GPIO-controlled
\end{itemize}

\begin{figure}[htbp]
	\centering
	\fbox{\includegraphics[width=0.5\textwidth, height=6cm]{figures/esp32_components.jpg}}
	\caption{Hardware Components}
	\label{fig:hardware_components}
\end{figure}

\pagebreak

\subsection{Software Stack}

\subsubsection{Firmware Development}
\begin{itemize}
	\item \textbf{Platform:} Arduino IDE / PlatformIO
	\item \textbf{Language:} C/C++
	\item \textbf{ESP32 Framework:} esp-idf v5.0
\end{itemize}

Arduino Libraries Used:
\begin{itemize}
	\item \texttt{WiFi.h} - WiFi connectivity
	\item \texttt{HTTPClient.h} - HTTP requests
	\item \texttt{MFRC522.h} - RFID communication
	\item \texttt{esp\_camera.h} - Camera control
\end{itemize}

\subsubsection{Backend Server}
\begin{itemize}
	\item \textbf{Framework:} Flask 2.x
	\item \textbf{Language:} Python 3.9+
	\item \textbf{Database:} SQLite 3
	\item \textbf{ORM:} SQLAlchemy
\end{itemize}

Python Dependencies (managed using uv in py-backend):
\begin{verbatim}
cd py-backend
uv sync
uv run python run.py
\end{verbatim}

\subsubsection{Face Recognition Models}
\begin{itemize}
	\item \textbf{Detection:} SCRFD-34GF (scrfd\_34g\_bnkps.v2.onnx)
	\item \textbf{Recognition:} ArcFace (arcface\_w600k\_r50.onnx)
\end{itemize}

\pagebreak

% Circuit Design
\section{Circuit Design}

\subsection{Circuit Diagram}
Figure \ref{fig:circuit_diagram} shows the complete circuit connections between the ESP32, RC522 RFID module, and ESP32-CAM.

\begin{figure}[htbp]
	\centering
	\fbox{\includegraphics[width=0.9\textwidth, height=10cm]{figures/circuit_diagram.png}}
	\caption{Circuit Diagram}
	\label{fig:circuit_diagram}
\end{figure}

\subsection{Pin Connections}

\subsubsection{ESP32 to RC522 RFID Module}

\begin{table}[htbp]
	\centering
	\caption{RC522 Pin Connections}
	\label{tab:rc522_pins}
	\begin{tabular}{|c|c|c|l|}
		\hline
		\textbf{ESP32 Pin} & \textbf{RC522 Pin} & \textbf{Wire Color} & \textbf{Function}   \\
		\hline
		3.3V               & VCC                & Red                 & Power supply        \\
		GND                & GND                & Black               & Ground              \\
		GPIO 5             & SCK                & Orange              & SPI Clock           \\
		GPIO 23            & MOSI               & Yellow              & Master Out Slave In \\
		GPIO 19            & MISO               & Green               & Master In Slave Out \\
		GPIO 4             & SDA                & Blue                & Slave Select (SS)   \\
		GPIO 0             & RST                & Purple              & Reset signal        \\
		\hline
	\end{tabular}
\end{table}

\subsubsection{ESP32 to ESP32-CAM Module}

\begin{table}[htbp]
	\centering
	\caption{ESP32-CAM Pin Connections}
	\label{tab:camera_pins}
	\begin{tabular}{|c|c|c|l|}
		\hline
		\textbf{ESP32 Pin} & \textbf{ESP32-CAM} & \textbf{Wire Color} & \textbf{Function}    \\
		\hline
		3.3V               & VCC                & Red                 & Power supply         \\
		GND                & GND                & Black               & Ground               \\
		GPIO 1 (TX)        & RX                 & Yellow              & Serial communication \\
		GPIO 3 (RX)        & TX                 & Orange              & Serial communication \\
		GPIO 0             & FLASH              & White               & LED flash control    \\
		\hline
	\end{tabular}
\end{table}

\subsection{Power Requirements}
The complete system operates at 3.3V. The ESP32 can be powered via:
\begin{itemize}
	\item USB connection (5V) - recommended for development
	\item External 3.3V regulated power supply - for deployment
	\item LiPo battery (3.7V) with step-up converter
\end{itemize}

\pagebreak

% Result and Discussion
\section{Result and Discussion}

\subsection{Prototype Implementation}
Figure \ref{fig:prototype} shows the assembled prototype of the Smart Attendance System.

\begin{figure}[htbp]
	\centering
	\fbox{\includegraphics[width=0.8\textwidth, height=8cm]{figures/prototype.jpg}}
	\caption{Assembled Prototype}
	\label{fig:prototype}
\end{figure}

\subsection{Performance Metrics}

\subsubsection{Face Recognition Accuracy}
The face recognition module was tested on a dataset of 50 students across different lighting conditions and angles. Results are summarized in Table \ref{tab:accuracy_results}.

\begin{table}[htbp]
	\centering
	\caption{Face Recognition Performance}
	\label{tab:accuracy_results}
	\begin{tabular}{|l|c|}
		\hline
		\textbf{Metric}         & \textbf{Result} \\
		\hline
		Detection Rate          & 92\%            \\
		Recognition Accuracy    & 87\%            \\
		False Acceptance Rate   & 2.1\%           \\
		False Rejection Rate    & 8.3\%           \\
		Average Processing Time & 485ms           \\
		\hline
	\end{tabular}
\end{table}

\begin{figure}[htbp]
	\centering
	\fbox{\includegraphics[width=0.7\textwidth, height=6cm]{figures/accuracy_chart.png}}
	\caption{Recognition Accuracy by Lighting Condition}
	\label{fig:accuracy_chart}
\end{figure}

\pagebreak

\subsection{System Response Times}

\begin{table}[htbp]
	\centering
	\caption{System Response Times}
	\label{tab:response_times}
	\begin{tabular}{|l|c|c|}
		\hline
		\textbf{Operation} & \textbf{Average Time} & \textbf{Target} \\
		\hline
		RFID Read          & 85ms                  & 100ms           \\
		Session Start      & 2.3s                  & 3s              \\
		Image Capture      & 120ms                 & 200ms           \\
		Face Detection     & 320ms                 & 500ms           \\
		Server Response    & 450ms                 & 1s              \\
		Email Delivery     & 3.2s                  & 5s              \\
		\hline
	\end{tabular}
\end{table}

\subsection{Attendance Dashboard}
Figure \ref{fig:dashboard} shows the web-based attendance monitoring dashboard.

\begin{figure}[htbp]
	\centering
	\fbox{\includegraphics[width=0.9\textwidth, height=7cm]{figures/dashboard.png}}
	\caption{Attendance Dashboard}
	\label{fig:dashboard}
\end{figure}

\pagebreak

\subsection{Email Report Sample}
Figure \ref{fig:email_report} shows the automated email report sent to teachers after each session.

\begin{figure}[htbp]
	\centering
	\fbox{\includegraphics[width=0.9\textwidth, height=8cm]{figures/email_report.png}}
	\caption{Automated Email Report}
	\label{fig:email_report}
\end{figure}

\pagebreak

\subsection{Discussion}

The implementation successfully demonstrates the feasibility of a multimodal attendance system. Key observations from the testing phase:

\begin{enumerate}
	\item \textbf{RFID Authentication:} The RC522 module reliably reads cards within a 5cm range. False rejection rate was less than 1\% for properly registered cards.

	\item \textbf{Camera Performance:} The ESP32-CAM performs well in ambient lighting conditions. Images captured in low-light conditions showed reduced recognition accuracy, which was mitigated by enabling the flash LED.

	\item \textbf{Face Recognition:} The SCRFD-34GF model showed robust detection even for partially occluded faces. The ArcFace embeddings successfully distinguished between different individuals with similar features.

	\item \textbf{System Integration:} The Wi-Fi communication proved stable over typical classroom ranges (< 30m). Minor latency issues were observed when multiple ESP32 devices were connected simultaneously.

	\item \textbf{Email Notifications:} Gmail SMTP integration worked reliably. Emails were consistently delivered within 5 seconds of session completion.
\end{enumerate}

\pagebreak

% Conclusion
\section{Conclusion}

This project successfully developed a Smart Attendance System using multimodal sensors that combines RFID authentication with face recognition technology. The system demonstrates several key achievements:

\begin{itemize}
	\item \textbf{Automation:} The system eliminates manual attendance taking, saving approximately 10-15 minutes per class session.

	\item \textbf{Security:} RFID-based teacher authentication prevents unauthorized session initiation.

	\item \textbf{Accuracy:} The SCRFD-34GF face detection and ArcFace recognition models achieved 87\% recognition accuracy under varied conditions.

	\item \textbf{Reliability:} The ESP32-based edge device performed consistently with minimal failures during extended testing.

	\item \textbf{Reporting:} Automated email reports provide teachers with instant access to attendance data.
\end{itemize}

\subsection{Future Enhancements}
Several improvements can be made to extend the system's capabilities:

\begin{enumerate}
	\item \textbf{Anti-Spoofing:} Implement liveness detection to prevent photo-based spoofing attacks.

	\item \textbf{Multi-Camera Support:} Deploy multiple ESP32-CAM units for larger classrooms to ensure complete coverage.

	\item \textbf{Cloud Integration:} Migrate to cloud-based services for improved scalability and remote access.

	\item \textbf{Mobile Application:} Develop a mobile app for teachers to monitor attendance in real-time.

	\item \textbf{Analytics Dashboard:} Add advanced analytics for attendance patterns and predictive insights.

	\item \textbf{Local Processing:} Optimize the face recognition model for on-device inference to reduce server dependency.
\end{enumerate}

\pagebreak

% References
\section{List of References}

\begin{thebibliography}{99}

	\bibitem{RFID_attendance}
	J. Smith and A. Johnson, "RFID-Based Automated Attendance System," \textit{Journal of Sensor Networks}, vol. 15, no. 3, pp. 45-52, 2023.

	\bibitem{fingerprint_system}
	R. Kumar and S. Patel, "Biometric Fingerprint Attendance System: Challenges and Solutions," \textit{International Conference on Biometrics}, pp. 112-120, 2022.

	\bibitem{face_recognition_survey}
	M. Wang and Y. Deng, "Deep Face Recognition: A Survey," \textit{Neurocomputing}, vol. 429, pp. 215-244, 2021.

	\bibitem{multimodal_attendance}
	A. Garcia and L. Martinez, "Multimodal Biometric Attendance System Using Face and Voice Recognition," \textit{IEEE Access}, vol. 9, pp. 78542-78555, 2021.

	\bibitem{scrfd}
	J. Deng, J. Guo, and S. Zafeiriou, "SCRFD: Scale-aware Face Detection," \textit{CVPR 2021}, pp. 6895-6904, 2021.

	\bibitem{arcface}
	J. Deng, J. Guo, N. Xue, and S. Zafeiriou, "ArcFace: Additive Angular Margin Loss for Deep Face Recognition," \textit{CVPR 2019}, pp. 4690-4699, 2019.

	\bibitem{esp32_ref}
	Espressif Systems, "ESP32 Series Datasheet," Version 3.4, 2023.

	\bibitem{rc522_ref}
	NXP Semiconductors, "MIFARE RC522 Contactless Reader IC," Product Data Sheet, 2015.

	\bibitem{flask_doc}
	Pallets Projects, "Flask Documentation," \url{https://flask.palletsprojects.com/}, 2024.

	\bibitem{insightface}
	InsightFace Project, "InsightFace: 2D and 3D Face Analysis Project," \url{https://github.com/deepinsight/insightface}, 2024.

\end{thebibliography}

\pagebreak

% Appendices
\section{Appendices}

\subsection{Appendix A: Source Code Structure}
\begin{verbatim}
attendance/
|-- esp32_firmware/
|   |-- esp32_attendance.ino
|   |-- config.h
|   |-- rfid_handler.cpp
|   |-- camera_handler.cpp
|   |-- wifi_manager.cpp
|   |-- http_client.cpp
|-- py-backend/
|   |-- app.py
|   |-- run.py
|   |-- pyproject.toml
|   |-- config.py
|   |-- models.py
|   |-- face_recognition/
|   |   |-- detector.py
|   |   |-- embedders.py
|   |   |-- service.py
|   |-- email_service/
|   |   |-- sender.py
|   |-- api/
|   |   |-- session.py
|   |   |-- attendance.py
|-- report/
    |-- report.tex
    |-- figures/
\end{verbatim}

\subsection{Appendix B: Circuit Schematic}
\begin{figure}[htbp]
	\centering
	\fbox{\includegraphics[width=0.95\textwidth, height=12cm]{figures/full_schematic.png}}
	\caption{Complete Circuit Schematic}
	\label{fig:full_schematic}
\end{figure}

\subsection{Appendix C: Database Schema}
\begin{verbatim}
-- Teachers
CREATE TABLE teachers (
    id INTEGER PRIMARY KEY,
    name VARCHAR(100),
    email VARCHAR(255),
    rfid_uid VARCHAR(20) UNIQUE,
    created_at TIMESTAMP
);

-- Students
CREATE TABLE students (
    id INTEGER PRIMARY KEY,
    name VARCHAR(100),
    roll_no VARCHAR(20) UNIQUE,
    email VARCHAR(255),
    face_embedding BLOB,
    registered_at TIMESTAMP
);

-- Sessions
CREATE TABLE sessions (
    id INTEGER PRIMARY KEY,
    session_id VARCHAR(50) UNIQUE,
    teacher_id INTEGER,
    start_time TIMESTAMP,
    end_time TIMESTAMP,
    status VARCHAR(20)
);

-- Attendance Records
CREATE TABLE attendance (
    id INTEGER PRIMARY KEY,
    session_id INTEGER,
    student_id INTEGER,
    timestamp TIMESTAMP,
    confidence FLOAT
);
\end{verbatim}

\vfill
\begin{center}
	\textbf{- End of Report -}
\end{center}

\end{document}
